\documentclass[10pt,twocolumn]{article}

\usepackage[margin=0.72in]{geometry}
\usepackage{amsmath,amssymb}
\usepackage{booktabs}
\usepackage{graphicx}
\usepackage{microtype}
\usepackage{xcolor}
\usepackage{hyperref}
\usepackage{cleveref}
\usepackage{enumitem}
\usepackage{multirow}
\usepackage{url}
\usepackage{balance}

\definecolor{tbdred}{RGB}{160,25,25}
\newcommand{\TBD}[1]{\textcolor{tbdred}{\textbf{[TBD: #1]}}}
\newcommand{\method}{\textsc{AMD-RTC}}
\newcommand{\pifive}{$\pi_{0.5}$}
\newcommand{\smolvla}{\textsc{SmolVLA}}

\hypersetup{
  colorlinks=true,
  linkcolor=blue,
  citecolor=blue,
  urlcolor=blue
}

\title{\method: Latency-Robust Vision--Language--Action Control\\
with Real-Time Chunking on AMD ROCm}

\author{
  Zhang Boyu\\
  AMD AI DevMaster Hackathon 2026 --- Track 3: Physical AI\\
  \texttt{boyuzhang1218@gmail.com} \\
  \texttt{https://github.com/zby1218}
}

\date{August 2026}

\begin{document}
\maketitle

\begin{abstract}
Large vision--language--action (VLA) policies improve semantic and
visuomotor generalization, but their inference latency conflicts with
high-frequency robot control.  Naive action-chunk execution either pauses the
robot while waiting for a new prediction or concatenates stale and newly
predicted chunks with discontinuities.  We present \method, an end-to-end
Physical AI pipeline that fine-tunes \pifive and a workshop \smolvla baseline on
synthetic Franka pick-and-place demonstrations and evaluates Sync, Async
queue, and Real-Time Chunking (RTC) under measured and injected inference
delay.  The complete pipeline---photorealistic Genesis data generation, BF16
policy adaptation, GPU rendering, policy inference, and paired
evaluation---runs on a single 48\,GB AMD Radeon GPU through PyTorch ROCm.
On Env-A, \smolvla Sync places $2/20$ at native delay but $0/20$ under
$+200$\,ms due to freeze-on-replan; Async recovers $1/20$ place and $4/20$
lift with zero freeze steps, while RTC guidance does not beat Async on place
success.  \pifive expert-only fine-tuning fails closed-loop grasp ($0/20$);
unfreezing vision (B1d) unlocks $1/20$ place.  A held-out Env-B runtime
fruit shift of $3$--$5$\,cm (paired seeds; never in training) zeros place
success for both policies ($0/20$).  Chunk inference P50 is approximately
$305$\,ms (\pifive) / $343$\,ms (\smolvla).
\end{abstract}

\section{Introduction}

Vision--language--action models map visual observations, proprioception, and
language instructions directly to robot actions.  Early systems such as
RT-2 demonstrated transfer from web-scale vision-language pretraining to
robot control \cite{zitkovich2023rt2}, while open models such as OpenVLA
made downstream adaptation accessible \cite{kim2024openvla}.  Recent
flow-matching VLAs, including $\pi_0$ and \pifive, generate continuous action
chunks suitable for dexterous control \cite{black2024pi0,black2025pi05}.

Action chunking reduces the effective decision horizon and improves temporal
coherence \cite{zhao2023act}.  It does not, however, remove the systems-level
mismatch between a high-latency foundation model and a control loop that must
advance at tens of Hertz.  Waiting for a complete chunk lowers throughput;
executing an old chunk while a new one is generated creates stale actions;
and directly switching chunks can introduce out-of-distribution command
jumps.  RTC addresses this issue at inference time by asynchronously
generating the next chunk and treating the overlap as a flow inpainting
problem \cite{black2025rtc}.

This work asks a deliberately narrow question: \emph{can RTC make a
fine-tuned VLA less sensitive to inference delay on an AMD Radeon Physical AI
stack?}  We avoid conflating this systems question with language
generalization, sim-to-real transfer, or reinforcement learning.  Our
planned artifacts are:

\begin{enumerate}[leftmargin=*,itemsep=2pt]
  \item an AMD ROCm pipeline for Genesis synthetic data, \pifive expert-only
        fine-tuning, and closed-loop Franka evaluation;
  \item a real-time Genesis adapter that advances control independently from
        policy inference and supports controlled delay injection;
  \item a causal ablation separating model scale, asynchronous execution, and
        RTC flow guidance; and
  \item a paired, fixed-manifest evaluation protocol with explicit uncertainty
        and complete per-timestep telemetry.
\end{enumerate}

\paragraph{Scope.}
The task is simulated single-object lifting.  We therefore evaluate latency
robustness and control continuity, not open-world semantics or real-world
transfer.  This boundary is essential: \pifive's broad pretraining motivates
its use, but does not by itself validate those capabilities in our setting.

\section{Related Work}

\paragraph{Vision--language--action policies.}
RT-2 formalized VLA co-training by representing robot actions as tokens
\cite{zitkovich2023rt2}.  OpenVLA released a 7B open-source policy trained on
970k robot trajectories and demonstrated parameter-efficient adaptation
\cite{kim2024openvla}.  $\pi_0$ introduced a VLM backbone with a flow-matching
action expert \cite{black2024pi0,lipman2023flow}; \pifive extended this family
through heterogeneous co-training and high-level semantic prediction
\cite{black2025pi05}.  We choose \pifive because RTC is defined for diffusion-
and flow-based policies and the model has an official LeRobot implementation
\cite{lerobot2026pi05docs}.

\paragraph{Efficient VLA baselines.}
\smolvla is a 450M-parameter VLA designed for single-GPU training and
consumer-grade deployment \cite{shukor2025smolvla}.  It is a stronger baseline
than a non-language behavior-cloning policy because it shares the VLA and
action-chunking paradigm while differing substantially in scale.  We do not
assume that the larger model must win: a simple single-task dataset may favor
the compact policy.

\paragraph{Action chunking and real-time execution.}
ACT predicts action sequences to reduce compounding error and introduced
temporal aggregation across overlapping predictions \cite{zhao2023act}.
RTC instead targets diffusion/flow policies at inference time.  It freezes
actions already committed for execution and inpaints the remainder of the new
chunk, permitting asynchronous prediction without discontinuous switches
\cite{black2025rtc}.  The original work evaluated dynamic simulation and
real-world bimanual tasks.  We study its integration into a photorealistic
Genesis--LeRobot--ROCm pipeline.

\paragraph{Open robot-learning infrastructure.}
LeRobot provides a unified stack for datasets, policies, training, and
asynchronous inference \cite{cadene2026lerobot}.  Genesis provides GPU-enabled
physics and rendering for embodied AI \cite{genesis2024}.  Their combination
allows the complete experiment to run in the competition's cloud environment
without physical hardware.

\section{Problem Formulation}

At control time $t$, a policy receives RGB observations
$I_t^{1:K}$, joint state $q_t$, and language instruction $\ell$, and predicts
an $H$-step action chunk
\begin{equation}
  A_t = \pi_\theta(I_t^{1:K}, q_t, \ell)
      = [a_t, a_{t+1}, \ldots, a_{t+H-1}].
\end{equation}
The control period is $\Delta t=1/30$\,s, while policy latency
$L$ may span multiple periods.  A synchronous executor requests $A_{t+H}$
only after consuming the current chunk and holds the previous command during
inference.  An asynchronous executor requests a future chunk early but may
switch discontinuously when it arrives.

RTC conditions the flow denoising process for the new chunk on the unexecuted
suffix of the previous chunk.  Conceptually, if $M$ marks committed or
overlapping steps, guidance penalizes disagreement:
\begin{equation}
  \mathcal{L}_{\mathrm{guide}}
  \propto \left\|M \odot
  \left(A_{\mathrm{new}}-A_{\mathrm{previous}}\right)\right\|_2^2.
\end{equation}
We use the official LeRobot RTC implementation rather than reimplementing its
sampler \cite{lerobot2026rtcdocs}.  Our contribution is the simulator control
adapter, delay instrumentation, and paired evaluation.

\section{System Design}

\subsection{Hardware and Software}

All stages run on a Radeon Cloud instance with a 48\,GB AMD Radeon GPU
(\texttt{gfx1100}, PCI \texttt{0x744b}), PyTorch 2.9.1 built for ROCm 7.2.1,
and Python 3.12 on Ubuntu 24.04.4.  \smolvla, \pifive, and RTC share the
workshop LeRobot 0.4.4 environment with a local transformers openpi
compatibility patch; we did not overwrite that stack with a separate
LeRobot 0.6.x install. Exact pins: \texttt{ENVIRONMENT.md}.

\begin{table}[t]
\centering
\caption{System configuration used for all formal scores.}
\label{tab:system}
\small
\begin{tabular}{ll}
\toprule
Component & Configuration \\
\midrule
GPU & AMD Radeon Graphics, gfx1100, 48\,GB \\
ROCm & 7.2.1 \\
PyTorch & 2.9.1+rocm7.2.1 \\
Simulator & Genesis 1.3.1 \\
Training stack & LeRobot 0.4.4 (+ openpi patch for $\pi_{0.5}$) \\
Precision & BF16 \\
OS & Ubuntu 24.04.4 LTS \\
\bottomrule
\end{tabular}
\end{table}

\subsection{Robot and Scene}

We use a simulated Franka Emika Panda: seven arm joints and two parallel
gripper joints.  Both state and action are 9-D joint positions.  Franka is a
standard research manipulator, its redundant 7-DoF workspace supports
top-down grasping, and a stable Genesis asset already exists.  Changing
embodiment would invalidate the collected demonstrations and confound the
execution study.

The scene is a photorealistic rustic kitchen mesh with a red cube on the
floor-origin workspace.  Physics runs at 30\,Hz with four substeps and PD
joint-position control.  A fixed overhead RGB camera provides global geometry;
a wrist-mounted RGB camera provides close-range grasp information.  Both
streams are $640\times480$.  A legacy dataset key calls the wrist stream
\texttt{observation.images.side}; the report distinguishes key name from
physical placement.  The generator records the instruction
\texttt{Pick up the cube.}, whereas the legacy evaluator defaults to
\texttt{Pick up the red cube.}; every formal run overrides that default with
the training string so language mismatch is not an uncontrolled variable.

\subsection{Real-Time Control Adapter}

The simulator/control worker advances every 33.3\,ms and never blocks on model
inference.  A policy worker receives the latest observation and writes action
chunks to a thread-safe queue.  Artificial delay is applied only to the policy
worker.  This distinction is critical: sleeping while the entire simulator is
paused changes wall-clock time but does not expose the controller to stale
actions.

The policy predicts $H=50$ actions, while every executor uses the same
$E=10$-step closed-loop execution window.  This replaces the legacy evaluator,
which synchronously replans every simulator step
(\texttt{action\_horizon=1}) and therefore does not implement the chunk
semantics needed for RTC.  The policy chunk length, execution horizon, and
latency measured in control steps are logged separately.

We implement three executors:

\begin{itemize}[leftmargin=*,itemsep=1pt]
  \item \textbf{Sync}: request after chunk exhaustion; hold the last command
        while waiting.
  \item \textbf{Async-only}: request early and directly append/switch to the
        returned chunk, without flow guidance.
  \item \textbf{RTC}: asynchronously request and merge with the official RTC
        inpainting guidance.
\end{itemize}

The initial RTC configuration follows the official recommendation:
execution horizon 10, maximum guidance weight 10.0, and exponential prefix
attention.  We do not tune these values on the test manifest.  A common
0.15\,rad/control-step command limiter is applied after all executors as a
safety filter.  Because this filter can itself suppress discontinuities, we
measure boundary jumps both before and after filtering and measure physical
jerk from the resulting joint trajectory.

\section{Dataset}

We generate 120 successful scripted pick-and-place demonstrations in Genesis
using seed 42 (171 attempts; failed \texttt{not\_in\_basket} episodes are
skipped).  Episodes run at 30\,Hz with dual RGB and 9-D joint state/action.
Videos are AV1-encoded in LeRobot v3 format.  Dataset statistics include mean,
standard deviation, and quantiles required by \pifive normalization.  The
dataset path is frozen before training because the existing generator appends
episodes when a repository already exists.

\begin{table}[t]
\centering
\caption{Synthetic training dataset (Env-A).}
\label{tab:data}
\small
\begin{tabular}{ll}
\toprule
Property & Value \\
\midrule
Robot & Franka Panda (7 arm + 2 gripper) \\
Episodes / frames & 120 / $\approx$30k \\
Control rate & 30\,Hz \\
Observations & overhead RGB, wrist RGB, joints \\
Action & 9-D absolute joint position \\
Task string & Pick up the fruit and place it in the basket. \\
Expert kept / attempts & 120/171 ($\approx70.2\%$) \\
Task count & 1 \\
\bottomrule
\end{tabular}
\end{table}

\paragraph{Why synthetic demonstrations?}
The competition permits simulation, all activity must remain in the cloud
environment, and no physical robot is available.  Scripted Cartesian
waypoints provide consistent supervision without claiming human demonstration
diversity.  The generator success rate validates the expert only; it is not a
policy result.

\paragraph{Data audit.}
Before training, we verify shapes, finite values, joint limits, random video
decoding, task-string consistency, action-difference tails, quantile
statistics, and a metadata hash.  No evaluation episode is read from the
training dataset.

\section{Policies and Training}

\subsection{\smolvla Baseline}

The \smolvla baseline uses the existing 450M checkpoint and recipe: BF16
autocast, frozen vision encoder, and trainable action expert and state
projection.  It is trained for 4000 steps (batch 4, seed 2026, peak VRAM
2.30\,GB).  This baseline anchors performance to an affordable VLA designed
for single-GPU robotics \cite{shukor2025smolvla}.

\subsection{\pifive Adaptation}

We first adapt \pifive in expert-only mode (freeze vision/VLM; train action
expert and projections) for 3000 steps (B1), then continue expert-only (B1c).
Closed-loop evaluation showed zero grasps.  We therefore continue from B1c
with full-model training and unfrozen vision for 2000 steps (B1d; batch 1,
peak VRAM 32.3\,GB), which unlocked the first place success.  Training uses
BF16 and gradient checkpointing in an isolated newer LeRobot environment so
the workshop 0.4.4 stack remains intact.  Seeds are fixed.  Model-family
comparisons are single-seed system results rather than estimates of
training-run variance.

We do not include LoRA in the main study.  Parameter-efficient tuning is a
valid separate topic, but would add a second adaptation variable and does not
address action-execution latency.

\section{Experimental Protocol}

\subsection{Fixed Paired Manifest}

All policies are evaluated on the same 30 initial object positions and seeds.
The manifest is generated once and frozen.  It contains 10 center, 10 edge,
and 10 near-corner positions within the training workspace.  Stratification
prevents a small random sample from overrepresenting easy central grasps.  The
evaluation is in-distribution with respect to position; we do not label it
OOD.  Every summary records the manifest ID, train and evaluation seeds,
checkpoint hash, language string, camera layout, renderer, chunk length,
execution horizon, and command-filter configuration.  Formal comparisons use
AMD GPU rendering only; CPU llvmpipe is reserved for debugging because the
existing pipeline exhibits a substantial renderer-dependent success-rate
shift.

\subsection{Baselines and Ablations}

\begin{table*}[t]
\centering
\caption{Pre-registered experiment matrix. P0 rows are mandatory. All \pifive
rows use the identical fine-tuned checkpoint.}
\label{tab:matrix}
\small
\begin{tabular}{llllll}
\toprule
ID & Policy & Executor & Extra delay & Purpose & Priority \\
\midrule
B0 & \smolvla & Sync & Native & Workshop-grade compact VLA baseline & P0 \\
B1-N & \pifive & Sync & Native & Isolate model change & P0 \\
O-N & \pifive & RTC & Native & RTC cost/benefit at measured latency & P0 \\
B1-200 & \pifive & Sync & +200\,ms & Delay-sensitive baseline & P0 \\
O-200 & \pifive & RTC & +200\,ms & Primary latency-robustness comparison & P0 \\
B2-200 & \pifive & Async-only & +200\,ms & Queue vs.\ flow-guidance mechanism & P1 \\
O-100 & \pifive & RTC & +100\,ms & Delay-response curve & P1 \\
O-H & \pifive & RTC & Native; $H_e\in\{5,10,15\}$ & Horizon trade-off & P2 \\
\bottomrule
\end{tabular}
\end{table*}

B0$\rightarrow$B1 isolates the model family.  B1$\rightarrow$B2 isolates
overlap between prediction and execution.  B2$\rightarrow$Ours isolates RTC
guidance.  B1$\rightarrow$Ours measures the complete system effect.  Comparing
only B0 to Ours would be uninterpretable because both model and executor
change.

\subsection{Metrics}

\paragraph{Success.}
An episode succeeds if maximum cube lift is at least 2\,cm, the lift is
sustained for eight frames, and final lift is at least 1\,cm.  We report
success count/total and a Wilson 95\% interval \cite{wilson1927interval}.

\paragraph{Task throughput.}
We report successes per wall-clock minute and time-to-success.  Failed episodes
remain in the denominator and are treated as timeouts.

\paragraph{Boundary discontinuity.}
For chunk boundaries $\mathcal{B}$ and actions normalized by training-set
standard deviations,
\begin{equation}
  D_{\mathrm{boundary}} =
  \frac{1}{|\mathcal{B}|}\sum_{t\in\mathcal{B}}
  \left\|\tilde a_t-\tilde a_{t-1}\right\|_2.
\end{equation}

\paragraph{Measured joint jerk.}
For measured joint positions $q_t$,
\begin{equation}
  J_t = \frac{\|q_t-3q_{t-1}+3q_{t-2}-q_{t-3}\|_2}{\Delta t^3}.
\end{equation}
The main jerk result excludes gripper joints so that intentional opening and
closing does not dominate arm smoothness.

\paragraph{Systems performance.}
We report training samples/s, median step time after warm-up, peak allocated
VRAM, checkpoint size, and chunk inference latency P50/P95/P99.  A 100-step
BF16-vs-FP32 microbenchmark measures the benefit of AMD mixed-precision
execution without introducing FP32-trained policy variants.

\subsection{Statistics}

The same episode IDs are paired across methods.  We report the full $2\times2$
success contingency table and exact McNemar test for paired binary outcomes
\cite{mcnemar1947}.  Continuous paired differences use bootstrap 95\%
intervals and are summarized with medians and interquartile ranges.  With only
20--30 episodes, intervals will be wide; we emphasize effect sizes and
individual paired transitions rather than overclaiming significance.

\section{Results}

Frozen snapshot: 2026-08-05. Task string:
\texttt{Pick up the fruit and place it in the basket.} Seed~$99$,
\texttt{chunk\_size}$=50$, Sync-full $E=50$, Async/RTC $E=10$,
\texttt{max\_steps}$=400$, shared \texttt{smooth\_action(0.15)}, AMD GPU
render only. Full log: \texttt{artifacts/formal/RESULTS.md}.

\subsection{Capability}

\begin{table}[t]
\centering
\caption{Closed-loop capability at native delay (Sync-full $E=50$).}
\label{tab:success}
\small
\begin{tabular}{lccc}
\toprule
Method & Place & Lift & Infer P50 \\
\midrule
B0 \smolvla Sync & $2/20$ & $3/20$ & $343$\,ms \\
B1 \pifive Sync (expert-only) & $0/20$ & $0/20$ & $304$\,ms \\
B1d \pifive Sync (unfreeze vision) & $1/20$ & $1/20$ & $305$\,ms \\
\bottomrule
\end{tabular}
\end{table}

\paragraph{Model comparison.}
On this small homogeneous pick-and-place dataset, workshop \smolvla remains
the stronger native Sync place baseline ($2/20$ vs.\ B1d $1/20$).
Expert-only \pifive never grasped in closed loop; continuing from B1c with full
training and unfrozen vision (B1d) unlocked the first place success. We
therefore treat B0$\rightarrow$B1d as a systems adaptation result, not a
claim that larger VLAs automatically win on simple kitchen tasks.

\subsection{Latency Robustness}

\begin{table}[t]
\centering
\caption{\smolvla latency story (primary system evidence).}
\label{tab:latency}
\small
\begin{tabular}{llccc}
\toprule
Delay & Executor & Place & Lift & Freeze \\
\midrule
Native & Sync-full & $2/20$ & $3/20$ & --- \\
$+200$\,ms & Sync-full & $0/20$ & $2/20$ & freezes \\
$+200$\,ms & Async $E=10$ & $1/20$ & $4/20$ & $0$ \\
$+200$\,ms & RTC $E=10$ & $0/20$ & $1/20$ & $0$ \\
\bottomrule
\end{tabular}
\end{table}

\paragraph{Primary finding.}
Injected $+200$\,ms delay collapses Sync place success through freeze-on-replan.
Continuous action-queue execution (Async) recovers place/lift and eliminates
freeze steps. RTC also eliminates freezes but does \emph{not} improve place
success over Async on this static task.

\paragraph{Mechanism.}
Because Async $\ge$ RTC on place success, we attribute the main delay-robustness
lever to overlapping prediction and execution rather than flow guidance.
An exploratory \pifive B1d $+200$\,ms table ($n=10$) is consistent with Sync
still hitting a graspable seed while Async/RTC $E=10$ remain at $0/10$; we
prefer the \smolvla $n=20$ table for the systems claim.

\subsection{Env-B Held-out Runtime Perturbation}

Capability-only probe (Sync-full $E=50$, native delay): identical Env-A fruit
inits (\texttt{seed=99}), then an independent perturb RNG
(\texttt{seed+7919}) applies $|\Delta xy|\sim U[3,5]$\,cm at a random angle;
basket fixed; never present in training. Mean $|\Delta xy|=3.7$\,cm.

\begin{table}[t]
\centering
\caption{Env-A vs Env-B place/lift (paired seeds).}
\label{tab:envb}
\small
\begin{tabular}{lcccc}
\toprule
Policy & Env-A place & Env-B place & Env-A lift & Env-B lift \\
\midrule
\smolvla B0 & $2/20$ & $0/20$ & $3/20$ & $1/20$ \\
\pifive B1d & $1/20$ & $0/20$ & $1/20$ & $1/20$ \\
\bottomrule
\end{tabular}
\end{table}

Every Env-A place success fails under its paired Env-B shift. We treat this as
a negative spatial-generalization result (pose sensitivity on a small dataset),
not a robustness claim; the primary contribution remains latency / non-stall
continuity.

\subsection{AMD ROCm Performance}

\begin{table}[t]
\centering
\caption{Measured AMD ROCm training and inference performance.}
\label{tab:amd}
\small
\begin{tabular}{lcc}
\toprule
Metric & \smolvla & \pifive \\
\midrule
Train steps / wall & $4000$ / $\sim10$\,min & B1 $3000$/$\sim37$\,min; B1d $+2000$/$\sim28$\,min \\
Peak train VRAM & $2.30$\,GB & $11.29$\,GB (B1) $\rightarrow$ $32.3$\,GB (B1d) \\
Chunk latency P50 & $343$\,ms & $304$\,ms \\
Chunk latency P95 & $361$\,ms & $518$\,ms \\
Mean train step time & $0.141$\,s & $\approx0.66$\,s (B1d) \\
\bottomrule
\end{tabular}
\end{table}

All reported scores use the AMD GPU renderer. We do not claim device power
figures without a validated sensor methodology.

\section{Discussion}

\paragraph{What RTC can and cannot improve.}
RTC changes action generation and execution, not semantic perception or task
knowledge.  It can reduce idle time and boundary discontinuity, but cannot
recover from an incorrectly localized object or a policy that never learned
the grasp.  The success metric and continuity metrics therefore answer
different questions.

\paragraph{Why a static task remains informative.}
Static lifting is less delay-sensitive than dynamic interception.  This makes
it a conservative test of throughput and continuity but a weak test of
delay-induced failure.  A valid outcome may be equal final success with lower
time-to-success.  We deliberately avoid adding a moving target after training,
which would conflate RTC with a severe distribution shift.

\paragraph{Practical value.}
Cloud and edge robots experience variable inference and communication delay.
An inference-time method that reuses the same checkpoint can improve system
responsiveness without recollecting demonstrations.  The released adapter and
telemetry protocol can be reused for dynamic tasks and other flow-based VLAs.

\section{Limitations}

The study uses one simulated robot, one object geometry, one scene, one
language instruction, and at most 30 paired evaluation states.  Synthetic
expert trajectories are smooth and homogeneous.  The task is static, so it
does not reproduce the most delay-sensitive scenarios in the RTC paper.
No real-robot or sim-to-real conclusion is supported.  Sample size limits
statistical power.  Finally, comparison between \smolvla and \pifive includes
differences in model architecture and training stack and should be interpreted
as a system baseline, not a controlled scaling law.

\section{Open-Source and Reproducibility}

We release:

\begin{enumerate}[leftmargin=*,itemsep=1pt]
  \item the Genesis RTC/latency adapter;
  \item fixed evaluation manifests;
  \item per-timestep telemetry and statistics scripts;
  \item version-locked environment specifications;
  \item checkpoints or download instructions;
  \item raw result JSON, figures, and representative videos; and
  \item a Docker recipe where licensing and storage permit.
\end{enumerate}

The competition submission includes this technical report, complete source,
a step-by-step English README, and a 3--5 minute demo video, following the
official requirements \cite{amd2026hackathon,amd2026submission}.

\section{Conclusion}

We delivered a reproducible AMD ROCm protocol for latency-aware VLA control:
synthetic demos, \smolvla/\pifive fine-tuning, and paired Sync/Async/RTC
evaluation under injected delay.  The strongest delay result is architectural
continuity---Async recovers place/lift versus Sync freeze---rather than an
automatic RTC place-success boost.  \pifive required vision unfreeze to grasp at
all.  Together these results define an honest latency benchmark for future
flow-based VLAs on AMD.

\section*{Team Contributions}

Solo submission by \textbf{Zhang Boyu}
(\texttt{boyuzhang1218@gmail.com}).
Responsible for end-to-end Track~3 engineering: Env-A data generation,
\smolvla/\pifive training on AMD ROCm, Sync/Async/RTC latency evaluation,
Env-B held-out perturbation, formal artifacts, technical report, and demo
video (\url{https://youtu.be/rXQQ4s-CxEE}).

\balance
\bibliographystyle{plain}
\bibliography{references}

\appendix
\onecolumn

\section{Reproduction Checklist}

\begin{enumerate}
  \item Verify AMD Radeon device (\texttt{gfx1100}, 48\,GB), ROCm 7.2.1, and BF16.
  \item Use the pinned LeRobot 0.4.4 + Genesis 1.3.1 environment
        (\texttt{ENVIRONMENT.md} / \texttt{requirements-lock.txt}).
  \item Obtain the 120-episode Env-A dataset and base checkpoints
        (paths in \texttt{DATA\_AND\_WEIGHTS.md}).
  \item Reproduce or load the frozen \smolvla B0 and \pifive B1d checkpoints.
  \item Run Sync-full ($E{=}50$) and Async/RTC ($E{=}10$) evals with seed~$99$.
  \item Optionally run Env-B via \texttt{scripts/run\_env\_b\_matrix.sh}.
  \item Generate summaries from raw JSON under \texttt{artifacts/formal/};
        do not hand-edit result tables.
  \item Confirm the demo at \url{https://youtu.be/rXQQ4s-CxEE} and this PDF.
\end{enumerate}

\section{Artifact Schema}

Each episode record should contain the experiment ID, checkpoint hash,
manifest ID, episode ID, object pose, injected delay, policy request and
completion timestamps, control timestamps, chunk IDs, queue lengths,
commanded actions, measured joint positions, cube poses, success state,
software versions, and GPU identifier.  This schema permits recomputation of
all reported metrics from raw data.

\section{Pre-Registered Interpretation Rules}

\begin{itemize}
  \item If RTC improves throughput but not success, claim execution efficiency,
        not capability improvement.
  \item If RTC only helps at $+200$\,ms, claim delay robustness, not universal
        superiority.
  \item If Async-only matches RTC, identify overlap as the dominant mechanism.
  \item If \smolvla matches or beats \pifive, report that task/data simplicity
        favors the compact policy.
  \item If all policies fail, investigate feature mapping, normalization, and
        action semantics before discussing model capacity.
\end{itemize}

\end{document}
